\subsection{Evaluation\revdeldec{ Protocols}\revadddec{: Why Generic Generation Metrics Are Insufficient}}
\label{sec:evaluation}

\revdelblock{ \subsection{Why Generic Generation Metrics are Insufficient}
}Evaluation remains one of the most critical, yet under-standardized, aspects of financial time-series generation. In conventional supervised learning or computer vision tasks, performance is often adequately measured by generic metrics like Mean Squared Error (MSE), Fréchet Inception Distance (FID)\revadddec{~\citep{heusel2017gans}}, or visual realism. However, FTSG involves multiple competing objectives where visual similarity or pointwise accuracy does not equate to financial validity.

A synthetic sequence might achieve a low MSE or MAE against a ground-truth trajectory, yet entirely fail to preserve crucial dynamics such as volatility clustering or tail-risk behavior, \revdeldec{ rendering it useless }\revadddec{limiting its usefulness }for stress testing. Conversely, generic metrics cannot evaluate whether generated data improves a trading agent's performance in a downstream simulation. Financial data's noise, non-stationarity, and regime-dependent nature demand that evaluation protocols go beyond generic reconstruction. In this section, we systematically organize evaluation practices into five dimensions, culminating in a unified evaluation checklist.

\subsection{\revdeldec{ Fidelity-Oriented }Evaluation\revadddec{: Fidelity-Oriented Assessment}}
A primary objective of Financial Time-Series Generation (FTSG) is to approximate the underlying multivariate data-generating distribution closely. Fidelity-oriented evaluation rigorously quantifies the statistical divergence between the synthetic sequences $\hat{\mathbf{X}}$ and the empirical observations $\mathbf{X}$. This assessment spans several mathematical dimensions:

\revdelblock{ 

\markdel{\textbf{Distributional Similarity:} To evaluate }}\revadddec{Distributional similarity evaluates }whether the model captures the marginal distributions of returns or prices \revdeldec{ , }\revadddec{through }optimal transport and probabilistic divergences\revdeldec{ are employed}. The 1-Wasserstein distance\revadddec{~\citebox{\citep{villani2009optimal,arjovsky2017wgan}} }between the empirical distribution $\mathbb{P}_r$ and the generated distribution $\mathbb{P}_g$ measures the minimum cost required to transform one distribution into the other:
    \begin{equation}
        W_1(\mathbb{P}_r, \mathbb{P}_g) = \inf_{\gamma \in \Pi(\mathbb{P}_r, \mathbb{P}_g)} \mathbb{E}_{(x,y)\sim \gamma}[\|x-y\|]
    \end{equation}
    where $\Pi(\mathbb{P}_r, \mathbb{P}_g)$ is the set of all joint distributions with marginals $\mathbb{P}_r$ and $\mathbb{P}_g$. Additionally, the \revdeldec{ Maximum Mean Discrepancy }\revadddec{maximum mean discrepancy }(MMD)\revdeldec{ is utilized to evaluate }\revadddec{~\citep{gretton2012kernel} evaluates }distances in a \revdeldec{ Reproducing Kernel Hilbert Space }\revadddec{reproducing kernel Hilbert space }(RKHS)\revadddec{~\citep{aronszajn1950theory} }$\mathcal{H}$ with kernel $k$:
    \begin{equation}
        \text{MMD}^2(\mathbb{P}_r, \mathbb{P}_g) = \mathbb{E}[k(x, x')] + \mathbb{E}[k(y, y')] - 2\mathbb{E}[k(x, y)]
    \end{equation}
    where $x, x' \sim \mathbb{P}_r$ and $y, y' \sim \mathbb{P}_g$.

\revdeldec{ 
\textbf{Temporal Dependence:} Financial data is }\revadddec{For temporal dependence, financial data are }fundamentally sequential. The generator must capture auto-regressive dynamics and lead-lag relationships. This is mathematically verified by comparing the Autocorrelation Function (ACF) of the generated returns $\hat{r}_t$ with the empirical returns $r_t$. For a given lag $k$, the ACF is defined as:
    \begin{equation}
        \rho_k = \frac{\text{Cov}(r_t, r_{t-k})}{\text{Var}(r_t)}
    \end{equation}
The Mean Absolute Error (MAE) between the empirical ACF vector and the synthetic ACF vector serves as a standard metric for temporal fidelity.

\revdeldec{ 
\textbf{Cross-Sectional Dependence:} In multivariate generation}\revadddec{For cross-sectional dependence}, the covariance structure among assets \revadddec{in multivariate generation }dictates portfolio-level risk. Cross-sectional fidelity is evaluated by computing the Frobenius norm of the difference between the empirical correlation matrix $\Sigma_r$ and the synthetic correlation matrix $\Sigma_g$:
    \begin{equation}
        \mathcal{E}_{corr} = \|\Sigma_r - \Sigma_g\|_F = \sqrt{\sum_{i=1}^N \sum_{j=1}^N (\sigma_{r, ij} - \sigma_{g, ij})^2}
    \end{equation}

\revdeldec{ 
\textbf{Spectral and Frequency-Domain Similarity:} To ensure }\revadddec{Spectral and frequency-domain similarity checks }that cyclical macroeconomic patterns and the low signal-to-noise ratio are maintained \revdeldec{ , spectral evaluation applies the Fast Fourier Transform }\revadddec{by applying the fast Fourier transform }(FFT)\revadddec{~\citep{cooley1965algorithm} }to convert sequences into the frequency domain. The Power Spectral Density (PSD) of the synthetic data is then compared against the real data to penalize generators that produce over-smoothed, frequency-deficient signals.

\subsection{\revadddec{Evaluation: }Financial Validity \revdeldec{ Evaluation}\revadddec{Assessment}}
A distinctive \revdeldec{ , non-negotiable }requirement of FTSG is that generated sequences must adhere to established financial theories and principles. High distributional similarity does not guarantee the preservation of empirical ``stylized facts''—\revdeldec{ the universal, }\revadddec{recurrent }non-linear statistical regularities observed across \revdeldec{ global financial markets. Evaluation must mathematically verify }\revadddec{many financial markets~\citep{cont2001empirical}. Evaluation should therefore test }the presence of these facts:

\revdelblock{ 

\markdel{\textbf{Heavy Tails and Kurtosis:} Financial returns are notoriously }}\revadddec{Heavy tails and kurtosis are assessed because financial returns are commonly }non-Gaussian, exhibiting excess kurtosis and fat tails. To quantify the generator's ability to produce extreme tail events (e.g., black swan crashes), researchers measure the sample kurtosis $\kappa = \frac{\mathbb{E}[(r_t - \mu)^4]}{\sigma^4}$ and estimate the tail index $\alpha$ using the Hill estimator\revadddec{~\citep{hill1975simple}}. For the upper tail, ordering observations as $X_{(1)} \ge X_{(2)} \ge \dots \ge X_{(n)}$, the Hill estimator for a threshold $k$ is:
    \begin{equation}
        \hat{\alpha} = \left( \frac{1}{k} \sum_{i=1}^k \ln X_{(i)} - \ln X_{(k+1)} \right)^{-1}
    \end{equation}
The synthetic sequence must yield an $\hat{\alpha}$ closely matching that of historical data.

\revdeldec{ 
\textbf{Volatility Clustering:} A fundamental market property where }\revadddec{Volatility clustering, documented by Mandelbrot~\citep{mandelbrot1963variation} and systematized among stylized facts by Cont~\citep{cont2001empirical}, is assessed as a market property in which }large price variations are followed by large variations. A valid generator must replicate this slow-decaying memory in variance. This is evaluated by computing the autocorrelation of squared returns $r_t^2$ or absolute returns $|r_t|$ over extended lags, verifying that $\rho_k(|r_t|) > 0$ for large $k$.

\revdeldec{ 
\textbf{Leverage Effect:} This asymmetric phenomenon dictates that }\revadddec{The leverage effect~\citebox{\citep{black1976studies,cont2001empirical}} is assessed as the asymmetric phenomenon in which }negative returns increase future volatility more severely than positive returns of the same magnitude. It is rigorously quantified by the cross-correlation between past returns and future squared returns at lag $\tau > 0$:
    \begin{equation}
        L(\tau) = \text{Corr}(r_t, r_{t+\tau}^2)
    \end{equation}
A financially valid synthetic dataset must exhibit $L(\tau) < 0$ for small $\tau$.

\revdeldec{ 
\textbf{Tail Risk Preservation (VaR and ES):} Synthetic }\revadddec{Tail risk preservation, assessed through VaR and ES, is checked because synthetic }data used for institutional stress-testing must not systematically underestimate market risk. Generative fidelity is audited using Value-at-Risk (VaR) and Expected Shortfall (ES)\revadddec{~\citebox{\citep{acerbi2002coherence,acerbi2014expected}}}. For a confidence level $1 - \alpha$ (e.g., $99\%$), VaR is the quantile of the return distribution:
    \begin{equation}
        \text{VaR}_\alpha = \inf \{ z \in \mathbb{R} \mid \mathbb{P}(r_t \le z) \ge 1 - \alpha \}
    \end{equation}
While ES measures the expected loss conditional on exceeding the VaR threshold:
    \begin{equation}
        \text{ES}_\alpha = \mathbb{E}[r_t \mid r_t \le \text{VaR}_\alpha]
    \end{equation}
Significant deviations between empirical and synthetic VaR/ES indicate a dangerous smoothing of downside risk.

\subsection{\revadddec{Evaluation: }Downstream Utility \revdeldec{ Evaluation}\revadddec{Assessment}}
To bridge the gap between theoretical machine learning metrics and practical finance, an increasingly adopted paradigm evaluates generated data \revdeldec{ strictly based on its }\revadddec{primarily by }utility in executing downstream financial tasks. \revdeldec{ If synthetic data enhances the performance of algorithmic systems, it possesses intrinsic }\revadddec{Improved downstream performance provides task-specific evidence of synthetic-data }value.

\revdelblock{ 

\markdel{\textbf{Forecasting Improvement (TSTR vs. TRTR):} This protocol }}\revadddec{Forecasting improvement, assessed through TSTR versus TRTR, }evaluates the generator as a data augmentation engine. A predictive model $f_\theta$ is trained entirely on generated data $\hat{\mathbf{X}}$ and evaluated on real, out-of-sample data $\mathbf{X}_{test}$. This \revdeldec{ Train-on-Synthetic, Test-on-Real }\revadddec{train-on-synthetic/test-on-real }(TSTR) performance\revadddec{, following the evaluation design of Esteban et al.~\citep{esteban2017real}, }is compared against the \revdeldec{ Train-on-Real, Test-on-Real }\revadddec{train-on-real/test-on-real }(TRTR) baseline. A TSTR predictive error (e.g., MSE or directional accuracy) that closely approaches or surpasses TRTR demonstrates exceptional generative utility.

\revdeldec{ 
\textbf{Trading Agent Training and Simulation:} By }\revadddec{Trading-agent training and simulation is evaluated by }deploying the generative model as an interactive market environment \revdeldec{ , Reinforcement Learning }\revadddec{in which reinforcement learning }(RL) trading agents can be trained on millions of synthetic episodes. The utility of the generator is subsequently quantified by the financial performance of the agent when deployed on real out-of-sample data, measured via the annualized Sharpe \revdeldec{ Ratio }\revadddec{ratio }($SR$)\revdeldec{ and Maximum Drawdown }\revadddec{~\citep{sharpe1994sharpe} and maximum drawdown }($MDD$):
    \begin{equation}
        SR = \frac{\mathbb{E}[R_p - R_f]}{\sqrt{\text{Var}(R_p)}}, \quad MDD = \max_{\tau \in (0, T)} \left( \frac{\max_{t \in (0, \tau)} P_t - P_\tau}{\max_{t \in (0, \tau)} P_t} \right)
    \end{equation}
    where $R_p$ is portfolio return, $R_f$ is the risk-free rate, and $P_t$ is the portfolio value at time $t$.

\subsection{\revadddec{Evaluation: }Robustness, Controllability, and Privacy}
As generative architectures—particularly foundation models and diffusion models—become increasingly sophisticated, evaluation protocols must expand to audit advanced capabilities and regulatory compliance.

\revdelblock{ 

\markdel{\textbf{Regime Robustness:} Generic }}\revadddec{Regime robustness is assessed because generic }metrics calculated over long, blended time horizons obscure a model's failure in specific contexts. Robustness evaluation partitions test data into discrete regimes (e.g., low-volatility bull markets vs. high-volatility crashes) using \revdeldec{ Hidden Markov Models (HMM)}\revadddec{hidden Markov models (HMMs)~\citebox{\citep{rabiner1989tutorial,hamilton1989new}} }or variance thresholds. Models are evaluated on their ability to maintain stable generative fidelity across these shifting boundaries.

\revdeldec{ 
\textbf{Conditional Consistency:} For }\revadddec{Conditional consistency is assessed for }multimodal or controllable generation (e.g., text-to-market or target-volatility synthesis)\revdeldec{ , evaluation }\revadddec{. Evaluation }must quantify semantic alignment. This involves measuring the correlation between the conditioning input (e.g., a target VaR level) and the actualized metric in the generated sequence, ensuring the generator is not merely ignoring the prompt to produce unconditional mode-collapsed outputs.

\revdeldec{ 
\textbf{Privacy Leakage and Memorization Risk:} A }\revadddec{Privacy leakage and memorization risk are assessed because a }critical motivation for FTSG is \revdeldec{ bypassing strict }\revadddec{reducing privacy risk under }financial data-sharing regulations (e.g., GDPR). To verify that a model has not simply memorized proprietary training data, the \revdeldec{ Distance-to-Closest Record }\revadddec{distance-to-closest-record }(DCR) \revadddec{statistic, used in privacy-preserving synthetic-data evaluation~\citep{bindschaedler2017plausible}, }is computed:
    \begin{equation}
        \text{DCR}(\hat{x}) = \min_{x \in \mathbf{X}_{train}} \|\hat{x} - x\|_2
    \end{equation}
If the DCR distribution peaks near zero, the model suffers from severe privacy leakage. Furthermore, models optimized with \revdeldec{ Differential Privacy }\revadddec{differential privacy }(DP)\revadddec{~\citebox{\citep{dwork2006differential,dwork2014algorithmic}} }are evaluated against their $(\epsilon, \delta)$-DP bounds, analyzing the mathematical trade-off between the privacy budget $\epsilon$ and the resulting degradation in downstream utility.

\subsection{\revadddec{Evaluation: }A Unified \revdeldec{ Evaluation }Checklist}
The lack of standardized evaluation frequently leads to contradictory claims regarding state-of-the-art performance in FTSG. To resolve this, we formalize a unified evaluation checklist, mapping the aforementioned rigorous metrics to their appropriate problem formulations.

As detailed in Table \ref{tab:evaluation_checklist}, researchers must construct a composite evaluation protocol. For instance, an \textit{imputation} model may prioritize pointwise accuracy (MSE) and temporal coherence, whereas a \textit{market simulation} architecture must be rigorously audited against stylized facts (\revdeldec{ Kurtosis, Ljung-Box testsfor volatility clustering}\revadddec{kurtosis and Ljung--Box tests~\citep{ljung1978measure} for temporal dependence}) and downstream utility (TSTR). No single metric is sufficient; financial validity emerges only at the intersection of distributional fidelity, economic logic, and practical applicability.

\subsection{\revadddec{Evaluation: }Emerging Trends \revdeldec{ in Evaluation }and Benchmarking}

The evaluation of synthetic financial data is rapidly transitioning from the isolated calculation of generic metrics toward comprehensive, standardized benchmark platforms. Recent literature \revdeldec{ reveals a definitive paradigm shift focused }\revadddec{indicates a shift in focus }on multimodal reasoning, structural market constraints, and full-pipeline financial workflows. We categorize these emerging trends into three primary directions:

The first direction is conditional specificity and scenario-based evaluation: standard distribution-matching metrics often fail to evaluate whether a conditional generative model actually respects its conditioning metadata (e.g., specific market regimes or macroeconomic textual prompts). To address this, \revadddec{Narasimhan and Jiang~}\citep{narasimhan2024timeweaver} introduced \textit{\revdeldec{ Time Weaver}\revadddec{TimeWeaver}}, \revdeldec{ pioneering novel }\revadddec{together with }specificity metrics that explicitly penalize generative models for failing to reproduce metadata-specific features in the synthesized sequences.

The second direction is market microstructure and LOB standardization: evaluating high-frequency tick data and Limit Order Books (LOB) requires specialized protocols due to their asynchronous nature and stringent cross-feature constraints. \revadddec{Zhong et al.~}\citep{lobench2025} recently introduced \textit{LOBench}, \revdeldec{ the first }\revadddec{a }standardized benchmark specifically engineered for LOB representation and generation. Utilizing massive, real-world China A-share market data, LOBench establishes unified preprocessing pipelines and task-specific baselines, allowing researchers to evaluate generative fidelity directly against microstructural properties such as queue dynamics and volume imbalances.

The third direction is full-pipeline workflows and multimodal reasoning: instead of treating time-series generation as an isolated algorithmic task, modern benchmarks evaluate synthetic data based on its utility within broader economic workflows. \revadddec{Jiang et al.~}\citep{finmaster2025} proposed \textit{FinMaster}, a holistic benchmark featuring a dedicated \textit{FinSim} module that evaluates how well generated synthetic, privacy-compliant financial data replicates market dynamics for downstream corporate auditing and accounting tasks. Similarly, \revdeldec{ comprehensive }\revadddec{Yu et al.~\citep{tsrbench2026} introduced the }multi-modal \revdeldec{ benchmarks like }\revadddec{benchmark }\textit{TSRBench} \revdeldec{ 
have been introduced }to stress-test the reasoning capabilities of Large Language Models (LLMs) over time-series data. These platforms reveal a critical insight for evaluation: strong semantic reasoning does not guarantee accurate context-aware numerical generation, highlighting the need for hybrid evaluation metrics that simultaneously score textual alignment and continuous mathematical fidelity.

